% Options for packages loaded elsewhere
\PassOptionsToPackage{unicode}{hyperref}
\PassOptionsToPackage{hyphens}{url}
\documentclass[
]{article}
\usepackage{xcolor}
\usepackage[margin=1in]{geometry}
\usepackage{amsmath,amssymb}
\setcounter{secnumdepth}{5}
\usepackage{iftex}
\ifPDFTeX
  \usepackage[T1]{fontenc}
  \usepackage[utf8]{inputenc}
  \usepackage{textcomp} % provide euro and other symbols
\else % if luatex or xetex
  \usepackage{unicode-math} % this also loads fontspec
  \defaultfontfeatures{Scale=MatchLowercase}
  \defaultfontfeatures[\rmfamily]{Ligatures=TeX,Scale=1}
\fi
\usepackage{lmodern}
\ifPDFTeX\else
  % xetex/luatex font selection
\fi
% Use upquote if available, for straight quotes in verbatim environments
\IfFileExists{upquote.sty}{\usepackage{upquote}}{}
\IfFileExists{microtype.sty}{% use microtype if available
  \usepackage[]{microtype}
  \UseMicrotypeSet[protrusion]{basicmath} % disable protrusion for tt fonts
}{}
\makeatletter
\@ifundefined{KOMAClassName}{% if non-KOMA class
  \IfFileExists{parskip.sty}{%
    \usepackage{parskip}
  }{% else
    \setlength{\parindent}{0pt}
    \setlength{\parskip}{6pt plus 2pt minus 1pt}}
}{% if KOMA class
  \KOMAoptions{parskip=half}}
\makeatother
\usepackage{color}
\usepackage{fancyvrb}
\newcommand{\VerbBar}{|}
\newcommand{\VERB}{\Verb[commandchars=\\\{\}]}
\DefineVerbatimEnvironment{Highlighting}{Verbatim}{commandchars=\\\{\}}
% Add ',fontsize=\small' for more characters per line
\usepackage{framed}
\definecolor{shadecolor}{RGB}{248,248,248}
\newenvironment{Shaded}{\begin{snugshade}}{\end{snugshade}}
\newcommand{\AlertTok}[1]{\textcolor[rgb]{0.94,0.16,0.16}{#1}}
\newcommand{\AnnotationTok}[1]{\textcolor[rgb]{0.56,0.35,0.01}{\textbf{\textit{#1}}}}
\newcommand{\AttributeTok}[1]{\textcolor[rgb]{0.13,0.29,0.53}{#1}}
\newcommand{\BaseNTok}[1]{\textcolor[rgb]{0.00,0.00,0.81}{#1}}
\newcommand{\BuiltInTok}[1]{#1}
\newcommand{\CharTok}[1]{\textcolor[rgb]{0.31,0.60,0.02}{#1}}
\newcommand{\CommentTok}[1]{\textcolor[rgb]{0.56,0.35,0.01}{\textit{#1}}}
\newcommand{\CommentVarTok}[1]{\textcolor[rgb]{0.56,0.35,0.01}{\textbf{\textit{#1}}}}
\newcommand{\ConstantTok}[1]{\textcolor[rgb]{0.56,0.35,0.01}{#1}}
\newcommand{\ControlFlowTok}[1]{\textcolor[rgb]{0.13,0.29,0.53}{\textbf{#1}}}
\newcommand{\DataTypeTok}[1]{\textcolor[rgb]{0.13,0.29,0.53}{#1}}
\newcommand{\DecValTok}[1]{\textcolor[rgb]{0.00,0.00,0.81}{#1}}
\newcommand{\DocumentationTok}[1]{\textcolor[rgb]{0.56,0.35,0.01}{\textbf{\textit{#1}}}}
\newcommand{\ErrorTok}[1]{\textcolor[rgb]{0.64,0.00,0.00}{\textbf{#1}}}
\newcommand{\ExtensionTok}[1]{#1}
\newcommand{\FloatTok}[1]{\textcolor[rgb]{0.00,0.00,0.81}{#1}}
\newcommand{\FunctionTok}[1]{\textcolor[rgb]{0.13,0.29,0.53}{\textbf{#1}}}
\newcommand{\ImportTok}[1]{#1}
\newcommand{\InformationTok}[1]{\textcolor[rgb]{0.56,0.35,0.01}{\textbf{\textit{#1}}}}
\newcommand{\KeywordTok}[1]{\textcolor[rgb]{0.13,0.29,0.53}{\textbf{#1}}}
\newcommand{\NormalTok}[1]{#1}
\newcommand{\OperatorTok}[1]{\textcolor[rgb]{0.81,0.36,0.00}{\textbf{#1}}}
\newcommand{\OtherTok}[1]{\textcolor[rgb]{0.56,0.35,0.01}{#1}}
\newcommand{\PreprocessorTok}[1]{\textcolor[rgb]{0.56,0.35,0.01}{\textit{#1}}}
\newcommand{\RegionMarkerTok}[1]{#1}
\newcommand{\SpecialCharTok}[1]{\textcolor[rgb]{0.81,0.36,0.00}{\textbf{#1}}}
\newcommand{\SpecialStringTok}[1]{\textcolor[rgb]{0.31,0.60,0.02}{#1}}
\newcommand{\StringTok}[1]{\textcolor[rgb]{0.31,0.60,0.02}{#1}}
\newcommand{\VariableTok}[1]{\textcolor[rgb]{0.00,0.00,0.00}{#1}}
\newcommand{\VerbatimStringTok}[1]{\textcolor[rgb]{0.31,0.60,0.02}{#1}}
\newcommand{\WarningTok}[1]{\textcolor[rgb]{0.56,0.35,0.01}{\textbf{\textit{#1}}}}
\usepackage{longtable,booktabs,array}
\usepackage{calc} % for calculating minipage widths
% Correct order of tables after \paragraph or \subparagraph
\usepackage{etoolbox}
\makeatletter
\patchcmd\longtable{\par}{\if@noskipsec\mbox{}\fi\par}{}{}
\makeatother
% Allow footnotes in longtable head/foot
\IfFileExists{footnotehyper.sty}{\usepackage{footnotehyper}}{\usepackage{footnote}}
\makesavenoteenv{longtable}
\usepackage{graphicx}
\makeatletter
\newsavebox\pandoc@box
\newcommand*\pandocbounded[1]{% scales image to fit in text height/width
  \sbox\pandoc@box{#1}%
  \Gscale@div\@tempa{\textheight}{\dimexpr\ht\pandoc@box+\dp\pandoc@box\relax}%
  \Gscale@div\@tempb{\linewidth}{\wd\pandoc@box}%
  \ifdim\@tempb\p@<\@tempa\p@\let\@tempa\@tempb\fi% select the smaller of both
  \ifdim\@tempa\p@<\p@\scalebox{\@tempa}{\usebox\pandoc@box}%
  \else\usebox{\pandoc@box}%
  \fi%
}
% Set default figure placement to htbp
\def\fps@figure{htbp}
\makeatother
\setlength{\emergencystretch}{3em} % prevent overfull lines
\providecommand{\tightlist}{%
  \setlength{\itemsep}{0pt}\setlength{\parskip}{0pt}}
\usepackage{bookmark}
\IfFileExists{xurl.sty}{\usepackage{xurl}}{} % add URL line breaks if available
\urlstyle{same}
\hypersetup{
  pdftitle={Tarea 02 - Gestión de Eventos y Participantes con Spring Boot: Refactorización progresiva},
  pdfauthor={José Manuel Sánchez Álvarez},
  hidelinks,
  pdfcreator={LaTeX via pandoc}}

\title{Tarea 02 - Gestión de Eventos y Participantes con Spring Boot:
Refactorización progresiva}
\author{José Manuel Sánchez Álvarez}
\date{}

\begin{document}
\maketitle

{
\setcounter{tocdepth}{2}
\tableofcontents
}
\subsection{\texorpdfstring{🧩 \textbf{Orden sugerido de
refactorización}}{🧩 Orden sugerido de refactorización}}\label{orden-sugerido-de-refactorizaciuxf3n}

\subsubsection{\texorpdfstring{\textbf{Fase 1 --- Controlador ``gordo''
(punto de partida, realizado en tarea
anterior)}}{Fase 1 --- Controlador ``gordo'' (punto de partida, realizado en tarea anterior)}}\label{fase-1-controlador-gordo-punto-de-partida-realizado-en-tarea-anterior}

El controlador hace todo:

\begin{itemize}
\tightlist
\item
  Accede al repositorio directamente.
\item
  Contiene la lógica de negocio y validaciones básicas.
\item
  Devuelve directamente \texttt{ResponseEntity} con entidades
  (\texttt{Evento}).
\end{itemize}

👉 Objetivo: que funcione el CRUD completo antes de modularizar.

\begin{center}\rule{0.5\linewidth}{0.5pt}\end{center}

\subsubsection{\texorpdfstring{\textbf{Fase 2 --- Separación de
responsabilidades}}{Fase 2 --- Separación de responsabilidades}}\label{fase-2-separaciuxf3n-de-responsabilidades}

Mover lógica de negocio al servicio.

\begin{enumerate}
\def\labelenumi{\arabic{enumi}.}
\item
  \textbf{Crear la capa de servicio} (\texttt{EventoService}):

  \begin{itemize}
  \tightlist
  \item
    Define los métodos \texttt{listarEventos},
    \texttt{obtenerEventoPorId}, \texttt{crearEvento},
    \texttt{actualizarEvento}, \texttt{borrarEvento}.
  \item
    Pasa el repositorio como dependencia (inyección).
  \end{itemize}
\item
  \textbf{Controlador → Servicio:}

  \begin{itemize}
  \tightlist
  \item
    El controlador \textbf{solo delega} en el servicio.
  \item
    Retorna el resultado de cada método del servicio.
  \item
    Ya no usa el repositorio directamente.
  \end{itemize}
\item
  \textbf{Objetivo didáctico:} comprender el principio
  \emph{``Separation of Concerns''} y el papel del servicio como
  intermediario.
\end{enumerate}

\begin{center}\rule{0.5\linewidth}{0.5pt}\end{center}

\subsubsection{\texorpdfstring{\textbf{Fase 3 --- Validaciones y
excepciones}}{Fase 3 --- Validaciones y excepciones}}\label{fase-3-validaciones-y-excepciones}

Refinar la capa de servicio.

\begin{enumerate}
\def\labelenumi{\arabic{enumi}.}
\tightlist
\item
  Añadir validaciones en el servicio (por ejemplo, comprobar si el
  evento existe).
\item
  Crear una excepción personalizada (\texttt{EventoNotFoundException}).
\item
  Crear un manejador global de excepciones (\texttt{@ControllerAdvice}).
\end{enumerate}

👉 Así el controlador queda limpio y coherente con buenas prácticas
REST.

\begin{center}\rule{0.5\linewidth}{0.5pt}\end{center}

\subsubsection{\texorpdfstring{\textbf{Fase 4 --- DTOs y
conversión}}{Fase 4 --- DTOs y conversión}}\label{fase-4-dtos-y-conversiuxf3n}

Evitar exponer entidades directamente.

\begin{enumerate}
\def\labelenumi{\arabic{enumi}.}
\tightlist
\item
  Crear \texttt{EventoDTO} o \texttt{EventoResponseDTO}.
\item
  Crear un \emph{mapper} o conversión sencilla (\texttt{ModelMapper} o
  manual).
\item
  Actualizar el servicio/controlador para trabajar con DTOs.
\end{enumerate}

👉 Esto prepara el terreno para incluir filtrado, paginación y HATEOAS
después.

\begin{center}\rule{0.5\linewidth}{0.5pt}\end{center}

\subsubsection{\texorpdfstring{\textbf{Fase 5 --- Paginación y
ordenación}}{Fase 5 --- Paginación y ordenación}}\label{fase-5-paginaciuxf3n-y-ordenaciuxf3n}

Integrar \texttt{Pageable}.

\begin{enumerate}
\def\labelenumi{\arabic{enumi}.}
\item
  Cambiar el método de listar:

\begin{Shaded}
\begin{Highlighting}[]
\AttributeTok{@GetMapping}
\KeywordTok{public}\NormalTok{ Page}\OperatorTok{\textless{}}\NormalTok{EventoDTO}\OperatorTok{\textgreater{}} \FunctionTok{listarEventos}\OperatorTok{(}\BuiltInTok{Pageable}\NormalTok{ pageable}\OperatorTok{)}
\end{Highlighting}
\end{Shaded}
\item
  En el servicio usar:

\begin{Shaded}
\begin{Highlighting}[]
\NormalTok{eventoRepository}\OperatorTok{.}\FunctionTok{findAll}\OperatorTok{(}\NormalTok{pageable}\OperatorTok{)}
\end{Highlighting}
\end{Shaded}
\item
  Probar URLs como:

\begin{verbatim}
/eventos?page=0&size=10&sort=fecha,asc
\end{verbatim}
\end{enumerate}

👉 Aquí se introducen conceptos de eficiencia y escalabilidad.

\begin{center}\rule{0.5\linewidth}{0.5pt}\end{center}

\subsubsection{\texorpdfstring{\textbf{Fase 6 --- Filtrado avanzado y
búsqueda}}{Fase 6 --- Filtrado avanzado y búsqueda}}\label{fase-6-filtrado-avanzado-y-buxfasqueda}

Añadir parámetros opcionales:

\begin{itemize}
\tightlist
\item
  \texttt{/eventos?nombre=Concierto}
\item
  \texttt{/eventos?fechaInicio=2025-05-10}
\end{itemize}

Puedes hacerlo con \texttt{@RequestParam} o
\texttt{Specification}/\texttt{Example}.

\begin{center}\rule{0.5\linewidth}{0.5pt}\end{center}

\subsubsection{\texorpdfstring{\textbf{Fase 7 --- HATEOAS y enlaces} (no
necesario, no
penaliza)}{Fase 7 --- HATEOAS y enlaces (no necesario, no penaliza)}}\label{fase-7-hateoas-y-enlaces-no-necesario-no-penaliza}

Añadir hipermedios si estás mostrando cómo diseñar una API REST
completa:

\begin{itemize}
\tightlist
\item
  Uso de \texttt{EntityModel}, \texttt{PagedModel}, \texttt{linkTo},
  \texttt{methodOn}.
\end{itemize}

\begin{center}\rule{0.5\linewidth}{0.5pt}\end{center}

\subsubsection{\texorpdfstring{\textbf{Fase 8 ---
Tests}}{Fase 8 --- Tests}}\label{fase-8-tests}

Cuando ya está modularizado:

\begin{itemize}
\tightlist
\item
  Unit tests al servicio con \texttt{@MockBean} o \texttt{Mockito}.
\item
  Tests de integración al controlador (\texttt{@SpringBootTest},
  \texttt{@AutoConfigureMockMvc}).
\end{itemize}

\begin{center}\rule{0.5\linewidth}{0.5pt}\end{center}

\subsection{\texorpdfstring{🧠 \textbf{Resumen
estructurado}}{🧠 Resumen estructurado}}\label{resumen-estructurado}

\begin{longtable}[]{@{}
  >{\raggedright\arraybackslash}p{(\linewidth - 4\tabcolsep) * \real{0.0625}}
  >{\raggedright\arraybackslash}p{(\linewidth - 4\tabcolsep) * \real{0.4062}}
  >{\raggedright\arraybackslash}p{(\linewidth - 4\tabcolsep) * \real{0.5312}}@{}}
\toprule\noalign{}
\begin{minipage}[b]{\linewidth}\raggedright
Fase
\end{minipage} & \begin{minipage}[b]{\linewidth}\raggedright
Enfoque
\end{minipage} & \begin{minipage}[b]{\linewidth}\raggedright
Objetivo didáctico
\end{minipage} \\
\midrule\noalign{}
\endhead
\bottomrule\noalign{}
\endlastfoot
1 & Controlador completo & Tener base funcional \\
2 & Crear capa de servicio & Separación de responsabilidades \\
3 & Validaciones y excepciones & Manejo correcto de errores \\
4 & DTOs & Buen diseño de API y encapsulación \\
5 & Paginación / Ordenación & Escalabilidad \\
6 & Filtrado & Consultas más útiles \\
7 & HATEOAS & API REST enriquecida (no necesario) \\
8 & Tests & Asegurar calidad y mantenimiento \\
\end{longtable}

\begin{center}\rule{0.5\linewidth}{0.5pt}\end{center}

\end{document}
